% Generated draft: requires mathematical and source review.
\documentclass[11pt]{article}
\usepackage[margin=25mm]{geometry}
\usepackage{fontspec}
\setmainfont{Latin Modern Roman}
\usepackage{amsmath,amssymb,mathtools,unicode-math}
\setmathfont{Latin Modern Math}
\usepackage{enumitem,needspace,xurl,hyperref}
\setlength{\parindent}{0pt}
\setlength{\parskip}{5pt}
\setlength{\emergencystretch}{3em}
\newcommand{\N}{\mathbb{N}}
\newcommand{\Z}{\mathbb{Z}}
\newcommand{\Q}{\mathbb{Q}}
\newcommand{\R}{\mathbb{R}}
\newcommand{\C}{\mathbb{C}}
\newcommand{\F}{\mathbb{F}}
\newcommand{\A}{\mathbb{A}}
\newcommand{\PP}{\mathbb{P}}
\newcommand{\E}{\mathbb{E}}
\newcommand{\eps}{\varepsilon}
\newcommand{\dd}{\,\mathrm{d}}
\begin{document}
\section*{UnsolvedMath ID: 30001090}
\textbf{Source problem code:} catalogue entry
\subsection{Definitions and mathematical statement}
Let $f:X\dashrightarrow X$ be a meromorphic self-map of a projective compact complex surface in the source's small-topological-degree setting. Conventionally, for a dominant map, its topological degree $\lambda_2(f)$ counts preimages of a generic point with multiplicity, its first dynamical degree $\lambda_1(f)$ measures exponential growth of pullbacks of divisor classes, and small topological degree means
\[
\lambda_2(f)<\lambda_1(f).
\]
The exact standing hypotheses, such as the model on which dynamical pullbacks and energy are considered, are not supplied.

Assume the report's finite-dynamical-energy hypotheses and let $\mu$ be the invariant probability measure constructed there. The original statement describes $\mu$ as a measure of maximal entropy. The question is whether it is the unique such measure. Writing $\mathcal P_f(X)$ for the admissible invariant Borel probability measures, the intended uniqueness statement is
\[
\forall\nu\in\mathcal P_f(X),\qquad
h_\nu(f)=h_\mu(f)=\sup_{\eta\in\mathcal P_f(X)}h_\eta(f)
\quad\Longrightarrow\quad \nu=\mu.
\]
Here $h_\nu(f)$ is measure-theoretic entropy. Because $f$ is meromorphic, invariance must be formulated on measures for which the map is defined almost everywhere and the pushforward is meaningful. The record does not define the class precisely. No finite-energy condition is silently imposed on every competing measure: finite dynamical energy is a hypothesis for the source construction of $\mu$, while the question asks uniqueness among maximal-entropy measures, not just among measures produced by that construction. Nor is the competitor restricted to ergodic or mixing measures.

Finite dynamical energy is a technical source condition whose definition and exact assumptions are absent. The source does not reproduce the measure construction or specify a topological-entropy convention for meromorphic maps, so no replacement energy criterion or numerical entropy identity is invented.

The corrected-verified note reports that the official report leaves uniqueness open under this energy hypothesis. The supplied triage checked 2026-08-21 labels the question open and says later papers establish maximal entropy and related structure without the general uniqueness claim. These are dated attributed statements, not new literature verification.
\subsection{Short English statement}
Under the report's finite-dynamical-energy hypotheses, ask whether its invariant measure for a small-topological-degree surface map is the only measure of maximal entropy. Competing measures are not silently restricted to the construction's finite-energy class.
\subsection{Sources}
\begin{itemize}
\item Local UnsolvedMath registry, permanent ID 30001090, source code catalogue entry; supplied complete record and corrected-verified report-check note. Literature checked 2026-08-21: the assessment is supplied dated curation, not independently verified here.
\item Trends and Developments in Complex Dynamics (2009), Oberwolfach Reports, DOI 10.4171/owr/2008/48, pp. 2703--2762: \url{https://doi.org/10.4171/owr/2008/48}. Proposed-year metadata gives 2008.
\item Diller--Dujardin--Guedj, small topological degree III: \url{https://arxiv.org/abs/0806.0146}.
\end{itemize}
\end{document}
